\documentclass[10pt]{article}

\usepackage[utf8]{inputenc}
\usepackage[english]{babel}
\usepackage[margin=0.85in]{geometry}
\usepackage{amsmath,amsthm,amssymb,stmaryrd}
\usepackage{hyperref,thmtools,enumitem,xcolor}

\hypersetup{colorlinks,linkcolor={red!50!black},citecolor={blue!50!black},urlcolor={blue!80!black}}

\declaretheorem[name=Theorem,numberwithin=section]{theorem}
\declaretheorem[name=Lemma,sibling=theorem]{lemma}
\declaretheorem[name=Proposition,sibling=theorem]{proposition}
\declaretheorem[name=Corollary,sibling=theorem]{corollary}
\declaretheorem[name=Definition,sibling=theorem,style=definition]{definition}
\declaretheorem[name=Remark,sibling=theorem,style=remark]{remark}
\declaretheorem[name=Conjecture,sibling=theorem]{conjecture}
\declaretheorem[name=Example,sibling=theorem,style=remark]{example}
\setlist[itemize]{parsep=0pt,topsep=2pt,before=\leavevmode}
\setlist[enumerate]{parsep=0pt,topsep=2pt,before=\leavevmode}
\setlength{\parindent}{1em}
\setlength{\parskip}{.15em}

\newcommand{\N}{\mathbb{N}}
\newcommand{\Prop}{\mathsf{Prop}}
\newcommand{\Type}{\mathsf{Type}}
\newcommand{\CIC}{\mathrm{CIC}}
\newcommand{\ZF}{\mathrm{ZF}}
\newcommand{\ZFC}{\mathrm{ZFC}}
\newcommand{\ZFh}{\mathrm{ZF}^{1/2}}
\newcommand{\EM}{\mathrm{EM}}
\newcommand{\Con}{\mathrm{Con}}
\newcommand{\rk}{\operatorname{rank}}
\newcommand{\dom}{\operatorname{dom}}
\newcommand{\PP}{\mathcal{P}}
\newcommand{\sem}[1]{\llbracket #1\rrbracket}
\newcommand{\hartogs}{\aleph}
\newcommand{\El}{\mathrm{El}}
\newcommand{\code}[1]{\ulcorner #1\urcorner}
\newcommand{\U}{\mathcal{U}}
\newcommand{\X}{\mathcal{X}}
\newcommand{\A}{\mathcal{A}}
\newcommand{\ind}{\mathsf{ind}}
\newcommand{\Acc}{\mathsf{Acc}}
\newcommand{\refl}{\mathsf{refl}}

\begin{document}

\title{Definable regularity:\\
  \large towards a set model of $\CIC$ with $\omega$ universes and excluded middle in $\ZF$}
\author{Mario Carneiro}
% Drafted with Claude; adjust authorship/acknowledgement as you see fit.
\date{Draft, \today}
\maketitle

\begin{abstract}
The Calculus of Inductive Constructions with an impredicative universe
of propositions, a tower of predicative universes, and excluded middle,
but without a choice or description operator, is usually modelled in
$\ZFC$ with one inaccessible cardinal per universe. We ask whether $\ZF$
alone proves its consistency, and we set out a model that would show
it. The function spaces of the model contain only the functions
definable by a term from parameters of bounded content; the extent
$\lambda_N$ of the universe $\Type_N$ is an ordinal with a regularity
property for such definable families only; and truth in the model,
which is not available while the extents are being determined, is
handled by evaluating terms relative to an arbitrary oracle for the
class of admitted elements and instantiating the oracle with the model
itself at the end. We prove that the model is sound provided the
extents exist, and we reduce their existence to a single statement, a
rank lemma saying that the types a term can denote below the extent of
a universe are bounded independently of that extent. The rank lemma is
open. This is a working paper: it fixes the type theory and the
definitions, records what is proved and what is not, and explains the
design decisions by the counterexamples that forced them.
\end{abstract}

\section{Introduction}

\paragraph{The question.}
Let $\CIC_\omega$ be the intensional type theory of
Section~\ref{sec:tt}: an impredicative, proof-irrelevant $\Prop$,
predicative cumulative universes $\Type_0:\Type_1:\cdots$, dependent
products and sums, $W$-types, natural numbers and finite types, an
identity type in $\Prop$ eliminating into every sort, propositional
truncation, quotients, an accessibility predicate eliminating into every
sort, propositional extensionality, and excluded middle in $\Prop$. It
has no choice operator, no description operator, and no unique choice.
Werner's set-theoretic model \cite{werner97} interprets $\Type_i$ as
$V_{\kappa_i}$ with $\kappa_i$ inaccessible, and with a description
operator the inaccessibles are necessary, since the sets-as-trees
interpretation then validates Replacement. Without one they are not
known to be. We conjecture:

\begin{conjecture}[main]\label{conj:main}
$\ZF\vdash\Con(\CIC_\omega+\EM)$.
\end{conjecture}

\paragraph{What is known.}
From below, $\CIC$ with $n+2$ universes and excluded middle interprets
$\ZFh_n$, Zermelo set theory with recursion terms and $n$ universes, and
proves its consistency \cite{paper2}; $\ZF$ proves the consistency of
all of these, so the lower bound is well below $\ZF$ and the boundary
is exactly Replacement for functional relations that are not given by
terms. From above, the inaccessibles cannot be removed from Werner's
model by replacing set-theoretic function spaces with realizability:
a set-sized universe of assemblies closed under sums over tracked
families has inaccessible size \cite[Prop.~4.4]{paper1}. The reason is
the context $X{:}\Type_0,\ f{:}X\to\Type_0$: whatever $f$ is allowed to
range over, $\Sigma x{:}X.\,f\,x$ must be in $\Type_0$, and the extent of
$\Type_0$ is a singular ordinal unless it is inaccessible, so some
small $X$ carries a cofinal family. No pointwise condition on $f$ can
exclude such a family (Example~\ref{ex:climb}). What excludes it in
$\ZFh$ is that a family there is a term, and the rank of a term's value
is bounded by an ordinal function of the ranks of its parameters, under
which the height of the model is closed \cite[\S3--4]{paper2}.

\paragraph{The model.}
We import that argument into type theory. The function spaces of the
model consist of the functions that are values of terms whose
parameters have bounded \emph{content}, the content of a parameter being
the supremum of the ranks of the small types occurring in it; and the
extent $\lambda_N$ of $\Type_N$ is \emph{$N$-closed}: as the parameters of
a term range over content below $\gamma<\lambda_N$, those values of the
term that are types of rank below $\lambda_N$ have rank below some
$\delta<\lambda_N$. This is what is left of the regularity of an
inaccessible when families are restricted to the definable ones, and
it makes the uniform half of type formation sound by definition. Two
difficulties remain.

The first is a circularity. The extents are determined from the bottom
up, and while $\lambda_N$ is being determined neither the higher
extents nor truth in the model are available: a proposition may
quantify over a function space, which is a set of definable functions,
whose definitions contain propositions. We resolve it by evaluating
terms relative to an \emph{oracle}, an arbitrary set $\X$ of admitted
elements: function spaces are cut down to $\X$, truth is computed from
the resulting decoding, and closedness quantifies over all oracles and
all higher extents. The generated elements $\A$ are the values of terms
under any oracle, and the model is the evaluation under the oracle
$\X:=\A$, which is one of those quantified over. No truth assignment
has to be guessed.

The second is the existence of $N$-closed ordinals. It follows from a
\emph{rank lemma} (Conjecture~\ref{conj:rank}): the types below the cut
that a term can denote are bounded by an ordinal that does not depend
on the cut. This is the analogue of the rank lemma of \cite{paper2} and
we believe it, having failed to refute it in several rounds; but its
proof is not an adaptation of Howard's majorizability argument
\cite{howard73}, for reasons given in Section~\ref{sec:ranklemma}, and
it is open. So the status of Conjecture~\ref{conj:main} is: reduced to
Conjecture~\ref{conj:rank}, modulo a soundness proof that is written by
cases and not yet checked rule by rule.

\paragraph{Why intensional.}
An earlier draft \cite[\S8]{paper2} worked with the extensional theory
of \cite{vdbergotten23,paper1}. Equality reflection makes the typing of
a term depend on the truth of the equations in its context, which
forced closedness to be stated for untyped terms, made every term
typable in some context, and defeated every attempt to stratify a
majorizability argument by types. The theory that proof assistants
implement is intensional, and for it none of this arises: typing is
decidable and independent of truth, and definitional equality is
computation, which holds in a set-theoretic evaluation outright. The
one place where a false equation could still do harm, the eliminator of
the identity type, is evaluated with a check.

\section{The type theory}\label{sec:tt}

\paragraph{Syntax.}
Sorts are $s::=\Prop\mid\Type_i$ for $i<\omega$. Terms:
\begin{align*}
t,u,A,B,C::={}&x\mid s\mid\Pi x{:}A.B\mid\lambda x{:}A.t\mid t\,u
  \mid\Sigma x{:}A.B\mid\langle t,u\rangle\mid\pi_1t\mid\pi_2t\\
 &\mid Wx{:}A.B\mid\mathsf{sup}\,t\,u\mid\ind^W_Ct
  \mid\mathbf n\mid\mathbf k_{\mathbf n}\mid\ind^{\mathbf n}_Ct_0\cdots t_{n-1}
  \mid\N\mid0\mid\mathsf S\,t\mid\ind^\N_Ct\,u\\
 &\mid t=_Au\mid\refl\,t\mid\ind^=_Ct
  \mid\|A\|\mid|t|\mid\ind^{\|\|}_Ct
  \mid A/R\mid[t]_R\mid\mathsf{lift}_Ct\,u\mid\mathsf{qind}_Pt\mid\mathsf{snd}_R\\
 &\mid\Acc_R\mid\mathsf{acc}\,t\,u\mid\ind^{\Acc}_Ct
  \mid\mathsf{propext}\mid\mathsf{em}
\end{align*}
with $\mathbf n$ the type with $n$ elements $\mathbf 0_{\mathbf n},\dots,(\mathbf{n{-}1})_{\mathbf n}$
for each $n<\omega$. We write $A\to B$ for $\Pi x{:}A.B$ with $x$ not
free in $B$, $\bot:=\|\mathbf 0\|$, $\neg P:=P\to\bot$, and
$P\vee Q:=\|\Sigma b{:}\mathbf 2.\,\ind^{\mathbf 2}_{\lambda\_.\Prop}\,P\,Q\,b\|$.

\paragraph{Judgments.}
$\vdash\Gamma$, $\Gamma\vdash t:A$, and $\Gamma\vdash t\equiv u:A$. We
write $\Gamma\vdash A:s$ for a type at a sort. The rules are those of a
pure type system with the following specifics.

\emph{Sorts and cumulativity.} $\Prop:\Type_0$ and $\Type_i:\Type_{i+1}$;
if $A:\Prop$ then $A:\Type_0$, and if $A:\Type_i$ then $A:\Type_{i+1}$.

\emph{Products.} If $A:s_1$ and $x{:}A\vdash B:s_2$ then
$\Pi x{:}A.B:\Prop$ when $s_2=\Prop$, and $\Pi x{:}A.B:\Type_{\max(i,j)}$
when $s_2=\Type_j$ and $s_1=\Type_i$, reading $\Prop$ as $\Type_0$ for
$s_1$. Abstraction and application as usual, with $\beta$ and $\eta$.

\emph{Sums, $W$-types, quotients.} If $A:\Type_i$ and
$x{:}A\vdash B:\Type_i$ then $\Sigma x{:}A.B:\Type_i$ and
$Wx{:}A.B:\Type_i$; if $A:\Type_i$ and $R:A\to A\to\Prop$ then
$A/R:\Type_i$. Pairs with projections, $\beta$ and $\eta$. Trees
$\mathsf{sup}\,a\,d:Wx{:}A.B$ for $a:A$ and $d:B[a/x]\to Wx{:}A.B$, with
$\ind^W_Cf:\Pi w{:}(Wx{:}A.B).\,C$ for a motive $w{:}W\vdash C:s$ in any
sort and
$f:\Pi a{:}A.\Pi d{:}(B[a/x]\to W).\,(\Pi b{:}B[a/x].\,C[d\,b/w])\to C[\mathsf{sup}\,a\,d/w]$,
and $\ind^W_Cf\,(\mathsf{sup}\,a\,d)\equiv f\,a\,d\,(\lambda b.\,\ind^W_Cf\,(d\,b))$.
Classes $[a]_R:A/R$ for $a:A$, the proof
$\mathsf{snd}_R:\Pi a\,a'{:}A.\,R\,a\,a'\to[a]_R=_{A/R}[a']_R$, and
$\mathsf{lift}_Cf\,h:A/R\to C$ for $C:s$ in any sort, $f:A\to C$, and
$h:\Pi a\,a'{:}A.\,R\,a\,a'\to f\,a=_Cf\,a'$, with
$\mathsf{lift}_Cf\,h\,[a]_R\equiv f\,a$; and the induction principle
$\mathsf{qind}_Pf:\Pi q{:}A/R.\,P\,q$ for $P:A/R\to\Prop$ and
$f:\Pi a{:}A.\,P\,[a]_R$.

\emph{Finite types and numbers.} $\mathbf n:\Type_0$ and $\N:\Type_0$ with
their constructors, and eliminators $\ind^{\mathbf n}_C\vec c:\Pi k{:}\mathbf n.\,C$
and $\ind^\N_Cc\,f:\Pi m{:}\N.\,C$ for motives in any sort, with the
usual computation rules.

\emph{Identity.} If $A:s$ and $a,a':A$ then $a=_Aa':\Prop$, with
$\refl\,a:a=_Aa$. For a motive $x{:}A,\,e{:}a=_Ax\vdash C:s$ in any sort
and $d:C[a/x,\refl\,a/e]$,
\[
\ind^=_Cd:\Pi x{:}A.\,\Pi e{:}a=_Ax.\,C,\qquad
\ind^=_Cd\,a\,(\refl\,a)\equiv d .
\]
There is no reflection rule.

\emph{Truncation.} If $A:s$ then $\|A\|:\Prop$, $|a|:\|A\|$ for $a:A$, and
$\ind^{\|\|}_Pf:\|A\|\to P$ for $P:\Prop$ and $f:A\to P$. It eliminates
into propositions only.

\emph{Accessibility.} For $X:s$ and $R:X\to X\to\Prop$,
$\Acc_R:X\to\Prop$, with
$\mathsf{acc}\,x\,h:\Acc_R\,x$ for $h:\Pi y{:}X.\,R\,y\,x\to\Acc_R\,y$. For a
motive $x{:}X\vdash C:s$ in any sort and
$f:\Pi x{:}X.\,(\Pi y{:}X.\,R\,y\,x\to C[y/x])\to C$,
\[
\ind^{\Acc}_Cf:\Pi x{:}X.\,\Acc_R\,x\to C,\qquad
\ind^{\Acc}_Cf\,x\,(\mathsf{acc}\,x\,h)\equiv f\,x\,(\lambda y\,r.\,\ind^{\Acc}_Cf\,y\,(h\,y\,r)).
\]

\emph{Proof irrelevance.} If $P:\Prop$ and $p,q:P$ then $p\equiv q:P$.

\emph{Axioms.} $\mathsf{propext}:\Pi P\,Q{:}\Prop.\,(P\to Q)\to(Q\to P)\to P=_\Prop Q$
and $\mathsf{em}:\Pi P{:}\Prop.\,P\vee\neg P$.

\emph{Conversion.} $\equiv$ is the congruence generated by the
computation rules above, $\beta$, $\eta$ for $\Pi$ and $\Sigma$, and
proof irrelevance; and if $t:A$ and $A\equiv B:s$ then $t:B$.

$\CIC_n$ is the theory with the sorts $\Prop,\Type_0,\dots,\Type_{n-1}$
only, and $\CIC_\omega$ is the union. Function extensionality is
derivable from the quotient rules by the usual argument. The theory is
close to the core of Lean~4 without $\mathsf{choice}$ \cite{carneiro19},
with $W$-types and $\Acc$ standing for the general inductive families,
and to the predicative calculus of inductive constructions with
$\mathsf{em}$, $\mathsf{propext}$ and quotients added. We have not
included general inductive families; we expect $W$, $=$ and $\Acc$ to
carry the features that matter for the strength, namely large
elimination over a well-founded structure of unbounded height and
singleton elimination.

\section{Evaluation relative to an oracle}\label{sec:sem}

All of this section is Werner's set-theoretic semantics \cite{werner97}
with three changes: universes have arbitrary ordinal extents; function
spaces, and the subtree functions of trees, are cut down to an arbitrary
set $\X$; and nothing is assumed to be sound. The evaluation is a total
$\ZF$-definable function, and what it is good for is decided afterwards.

\begin{definition}[structures, codes]\label{def:codes}
A \emph{structure} is an increasing sequence
$\vec\epsilon=(\epsilon_0<\epsilon_1<\cdots)$ of ordinals above $\omega$,
the \emph{extents}. We fix a representation of semantic objects as sets
in which natural numbers, the proof $\star$, pairs, functions, trees,
classes, and codes are pairwise distinct kinds, by tagging. The
\emph{codes} are: the constants $\bot,\top,\code\Prop,\code{\mathbf n},\code\N$,
of finite rank; the \emph{universe constants}
$\code{\U_i}:=\langle\mathsf U,\epsilon_i\rangle$, of rank about $\epsilon_i$;
and $\code\Pi(c,G)$, $\code\Sigma(c,G)$, $\code W(c,G)$, $\code/(c,R)$ for $c$
a code, $G$ a function into codes, and $R$ a set of pairs. The truth
values are the codes $\bot$ and $\top$. Put $U_i:=\{c:c\text{ a code},\ \rk(c)<\epsilon_i\}$,
so that $\code{\U_i}\in U_{i+1}\setminus U_i$.

Given $\lambda_0<\dots<\lambda_{N-1}$, an ordinal
$\alpha>\lambda_{N-1}$ (or $\alpha>\omega$ if $N=0$), and an increasing
sequence $\vec\alpha'$ of ordinals above $\alpha$, the
\emph{$(N,\alpha,\vec\alpha')$-structure} has extents
$\lambda_0,\dots,\lambda_{N-1},\alpha,\alpha'_{N+1},\alpha'_{N+2},\dots$:
the lower universes are given, $U_N$ is \emph{cut at} $\alpha$, the higher
universes are arbitrary.
\end{definition}

\begin{definition}[decoding relative to an oracle]\label{def:El}
An \emph{oracle} is a set $\X$. For a structure and an oracle, define
$\El^\X(c)$ by recursion on the rank of the code $c$:
\begin{itemize}
\item $\El^\X(\bot):=\emptyset$, $\El^\X(\top):=\{\star\}$,
  $\El^\X(\code\Prop):=\{\bot,\top\}$, $\El^\X(\code{\mathbf n}):=n$,
  $\El^\X(\code\N):=\omega$, $\El^\X(\code{\U_i}):=U_i$;
\item $\El^\X(\code\Pi(c,G))$ is the set of functions $f\in\X$ with
  $\dom f=\El^\X(c)$ and $f(a)\in\El^\X(G(a))$ for all $a$, if
  $\dom G\supseteq\El^\X(c)$, and $\emptyset$ otherwise;
\item $\El^\X(\code\Sigma(c,G))$ is the set of pairs $\langle a,b\rangle$
  with $a\in\El^\X(c)$ and $b\in\El^\X(G(a))$;
\item $\El^\X(\code W(c,G))$ is the set of well-founded trees, represented
  as sets of finite paths, whose nodes carry labels $a\in\El^\X(c)$ and
  have children indexed by $\El^\X(G(a))$, and such that at every node the
  function sending a child index to the subtree at that child is in $\X$;
\item $\El^\X(\code/(c,R))$ is the set of classes of $\El^\X(c)$ under the
  equivalence relation generated by $R\cap\El^\X(c)^2$.
\end{itemize}
By induction on codes $\El^\X(c)\subseteq V_{\rk(c)\cdot\omega}$; this is
why trees are sets of paths and not nested pairs.
\end{definition}

So a function of the decoding has \emph{exactly} the decoding of its
domain as its domain. Two consequences will matter. Functions are
extensional, so equality in the model can be equality of sets, and no
partial equivalence relation is needed. And the oracle is the only
unknown: once $\X$ is given, truth is computed.

\begin{definition}[evaluation]\label{def:eval}
Fix a structure and an oracle $\X$. Terms are \emph{sorted}: each type
former $\Pi$, $\Sigma$, $W$, $/$ that is not a proposition carries the
level $i$ of the sort $\Type_i$ at which it is formed, written
$\Pi^ix{:}A.B$ when it matters, and each subterm that is a proof, that is
whose type is a proposition, carries a mark; a derivation determines
the sorting and the marks, and typing is decidable, so this is
bookkeeping. For a sorted term $t$ and an environment $\rho$ assigning sets to variables
define $\sem t^\X\rho$ by recursion on $t$, returning $\emptyset$, which
is the code $\bot$, wherever a clause does not apply.
\begin{itemize}
\item $\sem x^\X\rho:=\rho(x)$; $\sem\Prop^\X\rho:=\code\Prop$,
  $\sem{\Type_i}^\X\rho:=\code{\U_i}$, $\sem{\mathbf n}^\X\rho:=\code{\mathbf n}$,
  $\sem\N^\X\rho:=\code\N$.
\item \emph{Type formers.} With $c:=\sem A^\X\rho$, the \emph{family at
  level $i$} of the binder $x{:}A$ with body $B$ is the function $G$ on
  $\El^\X(c)$ with $G(a):=\sem B^\X\rho[x{:=}a]$ if this is a code in $U_i$,
  and $G(a):=\bot$ otherwise. Then
  $\sem{\Pi^ix{:}A.B}^\X\rho:=\code\Pi(c,G)$, and likewise $\Sigma^i$ and
  $W^i$; and $\sem{A/R}^\X\rho:=\code/(c,R')$ with $R'$ the set of pairs
  $(a,a')\in\El^\X(c)^2$ with $\sem R^\X\rho(a)(a')=\top$. Families are
  \emph{sanitized}: their values lie in $U_i$ by fiat.
\item \emph{Propositions} are evaluated to truth values, computed from
  the decoding:
  $\sem{\Pi x{:}A.P}^\X\rho:=\top$ iff $\sem P^\X\rho[x{:=}a]=\top$ for
  all $a\in\El^\X(c)$, when $P$ is a proposition;
  $\sem{\|A\|}^\X\rho:=\top$ iff $\El^\X(c)\ne\emptyset$;
  $\sem{a=_Aa'}^\X\rho:=\top$ iff $\sem a^\X\rho=\sem{a'}^\X\rho$;
  $\sem{\Acc_R}^\X\rho$ is the function on $\El^\X(c)$, for $c$ the code of
  the carrier, that is $\top$ exactly on the accessible part of the
  relation $y\prec x:\iff\sem R^\X\rho(y)(x)=\top$ on $\El^\X(c)$.
  Every marked subterm, that is every proof, evaluates to $\star$; this
  covers $\refl$, $|{\cdot}|$, $\mathsf{acc}$, $\mathsf{snd}$,
  $\mathsf{qind}$, $\mathsf{propext}$, $\mathsf{em}$, and abstractions and
  applications of proofs.
\item $\sem{\lambda x{:}A.b}^\X\rho:=$ the function
  $a\mapsto\sem b^\X\rho[x{:=}a]$ on $\El^\X(c)$. It need not belong to
  $\X$. $\sem{t\,u}^\X\rho:=\sem t^\X\rho(\sem u^\X\rho)$ if the first is a
  function with the second in its domain, and $\emptyset$ otherwise.
\item Pairs and projections componentwise; $\mathsf{sup}\,a\,d$ is the
  tree with root label $\sem a^\X\rho$ and subtrees $\sem d^\X\rho(b)$;
  $[a]_R$ is the class of $\sem a^\X\rho$; numerals are numbers.
\item $\ind^{\mathbf n}$ evaluates by cases on the actual argument,
  $\ind^\N$ by primitive recursion, $\ind^W$ by recursion on the
  well-founded tree, all as functions on the decoding of the argument
  type. $\mathsf{lift}_Cf\,h$ is the function on the classes that sends
  $q$ to the common value of $\sem f^\X\rho$ on the members of $q$ if
  they agree, and to $\emptyset$ otherwise. $\ind^{\|\|}$ is constantly $\star$.
\item $\sem{\ind^=_Cd}^\X\rho$, where the motive is over $x{:}A$ and
  $e{:}a=_Ax$, is the function sending $x\in\El^\X(c)$ to the function
  that is $\star\mapsto\sem d^\X\rho$ if $x=\sem a^\X\rho$ and the empty
  function otherwise. So \emph{the eliminator of equality checks the
  equation}.
\item $\sem{\ind^{\Acc}_Cf}^\X\rho$ is defined by recursion along $\prec$
  on the accessible part $D$: at $x\in D$ it is
  $\star\mapsto\sem f^\X\rho(x)(\mathrm{rec}_x)$, where $\mathrm{rec}_x$ is
  the function on $\El^\X(c)$ with $\mathrm{rec}_x(y):=\{(\star,v_y)\}$, $v_y$
  the value at $y$, if $y\prec x$, and $\mathrm{rec}_x(y):=\emptyset$, the
  empty function, otherwise; at $x\notin D$ it is the empty function.
\end{itemize}
\end{definition}

\begin{remark}[what the evaluation does and does not assert]\label{rem:asserts}
The evaluation asserts nothing: not that a type lands in its universe,
nor that a term lands in the decoding of its type. For the oracle
$\X=V_\theta$ with $\theta$ large it is Werner's evaluation with
universes of arbitrary extent, and its failure to be a model is the
familiar one: $\sem{\Sigma x{:}X.\,f\,x}\rho$ is a code of rank
$\epsilon_0$ when $\rho(f)$ is cofinal, and with sanitized families a
pair $\langle a,b\rangle$ may fail to be in the decoding of its type
because the fibre was replaced by $\bot$. For a general oracle there is
a second failure: the value of an abstraction need not be in $\X$, and
then it is not an element of the function space and everything
downstream is junk. Soundness is a property of particular oracles, and
Section~\ref{sec:model} exhibits one.
\end{remark}

\section{Content, generated environments, closed ordinals}\label{sec:closed}

Throughout this section a level $N$ and ordinals
$\lambda_0<\dots<\lambda_{N-1}$ are fixed, the extents already
determined, and everything is relative to an
$(N,\alpha,\vec\alpha')$-structure and an oracle.

\begin{definition}[content]\label{def:content}
For a code $c$ and $e\in\El^\X(c)$ define the \emph{$N$-content}
$\|e\|_c$, an ordinal, by recursion on $c$:
\begin{itemize}
\item $\|e\|_c:=0$ for $c$ one of $\bot,\top,\code\Prop,\code{\mathbf n},\code\N$;
\item for $c=\code{\U_j}$ with any $j$, where $e$ is a code in $U_j$:
  $\|e\|_c:=\rk(e)+1$ if $e\in U_N$, that is if $\rk(e)$ is below the
  cut, and $\|e\|_c:=0$ otherwise;
\item $\|\langle a,b\rangle\|_{\code\Sigma(c',G)}:=\max(\|a\|_{c'},\|b\|_{G(a)})$;
\item $\|t\|_{\code W(c',G)}:=$ the supremum of $\|a\|_{c'}$ over the labels
  $a$ of $t$;
\item $\|q\|_{\code/(c',R)}:=\min\{\|a\|_{c'}:a\in q\}$;
\item $\|f\|_{\code\Pi(c',G)}:=\sup\{\|f(a)\|_{G(a)}:a\in\El^\X(c')\}$.
\end{itemize}
\end{definition}

The content of an element is the supremum of the ranks of the types of
level at most $N$ that occur in it as values, through all arguments and
components, fibre by fibre. The level of a code is read off the code and
not off the universe it is presented in, since $U_N\subseteq U_{N+1}$: a
family $\N\to\Type_{N+1}$ whose values lie in $U_N$ has the content of
those values. The identity on $U_N$ has content $\alpha$, the cut; so
does a family $\N\to U_N$ of cofinal ranks. If $c\in U_j$ with $j\le N$
then $\|e\|_c\le\rk(c)$ for every $e\in\El^\X(c)$, since a universe
constant mentioned by $c$ has extent below $\rk(c)$.

\begin{definition}[generated environments]\label{def:gen}
A \emph{generation tree} $\theta$ for a context $\Gamma$ assigns to each
variable $x$ of $\Gamma$ either the mark \emph{free} or a triple
$(\Gamma_x,t_x,\theta_x)$ of a context, a term typable in it, and a
generation tree for it; $\theta$ is required to be well founded, hence is
a finite tree, and there are countably many triples $(\Gamma,t,\theta)$
with $t$ typable in $\Gamma$. For an ordinal $\gamma$ and an oracle $\X$, a
\emph{$\gamma$-generated $\X$-environment for $(\Gamma,\theta)$} is an
environment $\rho$ such that, by recursion on $\theta$,
\begin{itemize}
\item $\rho$ is an \emph{$\X$-environment} for $\Gamma$:
  $\rho(x)\in\El^\X(\sem A^\X\rho)$ for every $x{:}A$ in $\Gamma$;
\item if $x$ is free then $\|\rho(x)\|_{\sem A^\X\rho}<\gamma$;
\item if $x$ has the triple $(\Gamma_x,t_x,\theta_x)$ then
  $\rho(x)=\sem{t_x}^{\X_x}\rho_x$ for some oracle $\X_x$ and some
  $\gamma$-generated $\X_x$-environment $\rho_x$ for $(\Gamma_x,\theta_x)$.
\end{itemize}
Each node may use its own oracle. Oracles may be taken to be subsets of
$V_\vartheta$ for $\vartheta$ the supremum of the extents, so for a fixed
structure the $\gamma$-generated environments form a set.
\end{definition}

\begin{definition}[closed ordinals]\label{def:closedord}
An ordinal $\lambda>\lambda_{N-1}$ (or $\lambda>\omega$ if $N=0$) is
\emph{$N$-closed} if $\beta\cdot\omega<\lambda$ for every $\beta<\lambda$,
and for every triple $(\Gamma,t,\theta)$ and every $\gamma<\lambda$ there
is $\delta<\lambda$ such that: for every sequence $\vec\alpha'$ of higher
extents, every oracle $\X$, and every $\gamma$-generated $\X$-environment
$\rho$ for $(\Gamma,\theta)$ in the $(N,\lambda,\vec\alpha')$-structure,
\[
\text{if }\sem t^\X\rho\text{ is a code of rank below }\lambda
\text{ then its rank is below }\delta.
\]
The extent $\lambda_N$ is the least $N$-closed ordinal, if there is one.
\end{definition}

Closedness is a \emph{gap condition}: a term whose parameters have
content below $\gamma$ cannot take, as the parameters, the oracle, and the
higher extents vary, values in $U_N$ of rank cofinal in $\lambda$. It is
stated in the structure cut at $\lambda$ itself and refers to no higher
extent. It bounds only the values that lie in $U_N$; whether the value
of a term of type $\Type_N$ does lie in $U_N$ is the business of the
soundness induction. For instance, for a definable operator
$t:\Type_N\to\Type_N$ the term $\ind^\N\,\N\,(\lambda\_.\,t)\,x$ with $x{:}\N$
free takes the values $t^n(\N)$; if these are all in $U_N$ they are
bounded below $\lambda$, and then $\Sigma n{:}\N.\,t^n(\N)$ is in $U_N$.
So an $N$-closed ordinal is not reached by iterating a definable type
operator, however fast that operator raises ranks
(Example~\ref{ex:sat}).

\begin{conjecture}[rank lemma]\label{conj:rank}
For every $N$, every triple $(\Gamma,t,\theta)$, and every ordinal
$\gamma$ there is an ordinal $G_{t,\theta}(\gamma)$ such that, for every cut
$\alpha>\max(\gamma,\lambda_{N-1})$, every $\vec\alpha'$, every oracle $\X$,
and every $\gamma$-generated $\X$-environment $\rho$ for $(\Gamma,\theta)$ in
the $(N,\alpha,\vec\alpha')$-structure: if $\sem t^\X\rho$ is a code of
rank below $\alpha$ then its rank is below $G_{t,\theta}(\gamma)$.
\end{conjecture}

\begin{lemma}[existence of closed ordinals]\label{lem:closedex}
Assume Conjecture~\ref{conj:rank} at level $N$. Then an $N$-closed
ordinal exists, so $\lambda_N$ is defined.
\end{lemma}
\begin{proof}
Take $G_{t,\theta}(\gamma)$ least, so that it is monotone in $\gamma$, and
enumerate the triples as $(\Gamma_k,t_k,\theta_k)$. Put
$\alpha_0:=\lambda_{N-1}+1$ (or $\omega+1$) and
$\alpha_{n+1}:=\max(\{\alpha_n\cdot\omega\}\cup\{G_{t_k,\theta_k}(\alpha_n):k\le n\})+1$,
and let $\lambda:=\sup_n\alpha_n$; then $\beta\cdot\omega<\lambda$ for
$\beta<\lambda$. Given the $k$-th triple and $\gamma<\lambda$, choose
$n\ge k$ with $\gamma\le\alpha_n$ and put
$\delta:=G_{t_k,\theta_k}(\alpha_n)<\alpha_{n+1}<\lambda$. The structure
cut at $\lambda$ is one of those the conjecture quantifies over, so in
it every value of $t_k$ at a $\gamma$-generated environment that is a
code of rank below $\lambda$ has rank below
$G_{t_k,\theta_k}(\gamma)\le\delta$.
\end{proof}

\section{The model}\label{sec:model}

Assume that $\lambda_N$ exists for every $N$, and work in the \emph{true
structure}, whose extents are the $\lambda_N$. For each $N$ it is the
$(N,\lambda_N,\vec\lambda_{>N})$-structure. Let
$\lambda_\omega:=\sup_N\lambda_N$; oracles are subsets of
$V_{\lambda_\omega}$.

\begin{definition}[generated elements]\label{def:A}
For each $N$ let $\A_N$ be the set of all values $\sem t^\X\rho$ in
$V_{\lambda_\omega}$, for a triple $(\Gamma,t,\theta)$, an oracle $\X$, and
a $\gamma$-generated $\X$-environment $\rho$ for $(\Gamma,\theta)$ with
$\gamma<\lambda_N$, the content being the $N$-content. The set of
\emph{generated elements} is $\A:=\bigcap_N\A_N$. The \emph{model} is the
evaluation relative to the oracle $\A$: its decoding is $\El:=\El^{\A}$,
its interpretation is $\sem t\rho:=\sem t^{\A}\rho$, and its environments
for $\Gamma$ are the $\A$-environments.
\end{definition}

The definition is not circular: $\A$ is a union over \emph{all} oracles,
and the oracle $\A$ is one of them. Truth in the model is whatever the
evaluation relative to $\A$ computes; it was never needed in order to
define $\A$.

\begin{lemma}[generated elements]\label{lem:gen}
(1) If $e\in\El^\X(c)$ for some oracle $\X$ and code $c$ in a universe,
and $\|e\|_c<\lambda_N$, then $e\in\A_N$. (2) If $\rho$ is an
$\X$-environment for $\Gamma$ with all values in $\A_N$ and $t$ is typable
in $\Gamma$, then $\sem t^\X\rho\in\A_N$. (3) Every code in a universe is
in $\A$; and if $c\in U_i$ and $e\in\El^\X(c)$ then $e\in\A_N$ for every
$N\ge i$. (4) $\El(c)\subseteq\A$ for every code $c$ in a universe.
\end{lemma}
\begin{proof}
(1) Take the context $X{:}\Type_i,\ x{:}X$ with both variables free,
$\rho(X):=c\in U_i$, $\rho(x):=e$, and $t:=x$. The content of $c$ as an
element of $\code{\U_i}$ is $\rk(c)+1$ if $c\in U_N$ and $0$ otherwise, in
either case below $\lambda_N$. (2) Graft the generation trees of the
values onto the variables of $\Gamma$; there are finitely many, so the
maximum of their bounds is below $\lambda_N$; and $\rho$ is an
$\X$-environment by hypothesis. (3) The first part was shown in (1); for
the second, $\|e\|_c\le\rk(c)<\lambda_i\le\lambda_N$. (4) By induction on
$c$. A function in $\El(\code\Pi(c,G))$ is in $\A$ by definition. A pair,
a tree, or a class of generated elements is the value of the
constructor term at a generated environment; for a tree the subtree
functions are in $\A$ by the definition of the decoding, and for a
class the relation enters as a variable $R:A\to A\to\Prop$ whose value
has content $0$. Codes are covered by (3), and numbers, truth values
and $\star$ by (1).
\end{proof}

So only the types that are not small at level $N$ are restricted by
$\A_N$: a function $\Type_N\to\Type_N$ must be the value of a term whose
parameters have content below $\lambda_N$, while every function between
two types in $U_N$ is its own parameter, so that the small types decode
as in Werner's model.

\begin{lemma}[type formation]\label{lem:tyform}
Let $\rho$ be an environment of the model for $\Gamma$, suppose
$\Gamma,x{:}A\vdash B:\Type_N$ is derivable, and suppose
$c:=\sem A\rho\in U_N$. Then $\sem{\Pi^Nx{:}A.B}\rho$,
$\sem{\Sigma^Nx{:}A.B}\rho$, and $\sem{W^Nx{:}A.B}\rho$ lie in $U_N$;
likewise for quotients.
\end{lemma}
\begin{proof}
The values of $\rho$ lie in $\A\subseteq\A_N$ by Lemma~\ref{lem:gen}(4),
so each is $\sem{t_y}^{\X_y}\rho_y$ for a $\gamma_y$-generated environment
with $\gamma_y<\lambda_N$. Let $\theta$ be the generation tree for
$\Gamma,x{:}A$ that gives each $y$ of $\Gamma$ its triple and marks $x$
free, and let $\gamma:=\max(\{\gamma_y\}\cup\{\rk(c)+1\})<\lambda_N$. For
every $a\in\El(c)$ we have $\|a\|_c\le\rk(c)<\gamma$, so $\rho[x{:=}a]$ is a
$\gamma$-generated $\A$-environment for $(\Gamma,x{:}A;\theta)$ in the true
structure. The family $G$ at level $N$ takes at $a$ the value
$\sem B\rho[x{:=}a]$ if this is a code of rank below $\lambda_N$, and
$\bot$ otherwise. Since $\lambda_N$ is $N$-closed there is
$\delta<\lambda_N$, depending only on $(B,\theta,\gamma)$, that bounds the
rank of every such value; so every value of $G$ has rank below $\delta$.
The code $\code Q(c,G)$ is a pair of $c$ and $G$, a function on
$\El(c)\subseteq V_{\rk(c)\cdot\omega}$; so its rank is below
$\max(\rk(c)\cdot\omega,\delta)+\omega<\lambda_N$.
\end{proof}

The division of labour is this. That an \emph{individual} fibre
$\sem B\rho[x{:=}a]$ lies in $U_N$ is the induction hypothesis of the
soundness proof for the premise $\Gamma,x{:}A\vdash B:\Type_N$; so in the
model the sanitization changes nothing. That the fibres are
\emph{uniformly} bounded, so that the code of the family lies in $U_N$,
is closedness.

\begin{lemma}[soundness]\label{lem:sound}
Let $\rho$ be an environment of the model for $\Gamma$.
\begin{enumerate}
\item If $\Gamma\vdash A:\Type_i$ then $\sem A\rho\in U_i$; if
  $\Gamma\vdash P:\Prop$ then $\sem P\rho$ is a truth value.
\item If $\Gamma\vdash t:A$ then $\sem t\rho\in\El(\sem A\rho)$.
\item If $\Gamma\vdash t\equiv u:A$ then $\sem t\rho=\sem u\rho$.
\end{enumerate}
\end{lemma}
\begin{proof}
By simultaneous induction on derivations, with the usual substitution
lemma $\sem{t[u/x]}\rho=\sem t\rho[x{:=}\sem u\rho]$. This is Werner's
argument; we go through the cases that differ.
\emph{Abstraction.} $\sem{\lambda x{:}A.b}\rho$ is a function on
$\El(\sem A\rho)$ with values in the right fibres by the induction
hypothesis for $b$ at the environments $\rho[x{:=}a]$, which are
environments of the model; and it is in $\A$ by
Lemma~\ref{lem:gen}(2),(4). So it is in the decoding of the product.
\emph{Application.} The argument lies in the decoding of the domain,
which is exactly the domain of the function.
\emph{Type formers} at sort $\Type_i$: at every $a\in\El(\sem A\rho)$ the
fibre $\sem B\rho[x{:=}a]$ lies in $U_i$ by the induction hypothesis for
$B$, so it is the value of the sanitized family there, and the code
lies in $U_i$ by Lemma~\ref{lem:tyform} at level $i$, the premises
$A:\Type_i$ and $B:\Type_i$ being available by cumulativity. Propositions
are truth values, in $U_0$, so the impredicative product needs no bound.
\emph{Trees.} The subtree function of $\mathsf{sup}\,a\,d$ at the root is
$\sem d\rho$, which is in $\A$, and at deeper nodes it is the subtree
function of a subtree, which is an element of the model.
\emph{Computation rules} hold as equations between sets: $\beta$ by the
substitution lemma; $\eta$ because a function of the decoding has
exactly the decoding of its domain as its domain; the rules for
$\ind^{\mathbf n}$, $\ind^\N$, $\ind^W$ by the recursions that define
them; $\ind^=_Cd\,a\,(\refl\,a)\equiv d$ because the check succeeds. For
$\ind^{\Acc}$ the right-hand side applies $\sem f\rho(x)$ to the value of
$\lambda y\,r.\,\ind^{\Acc}_Cf\,y\,(h\,y\,r)$, which is the function
$\mathrm{rec}_x$ of Definition~\ref{def:eval} on the nose, the inner
abstraction having the domain $\El$ of the truth value of $R\,y\,x$; and
$\mathrm{rec}_x$ is in $\A$ because it is the value of that term. For
$\mathsf{lift}_Cf\,h\,[a]_R\equiv f\,a$: the premise $h$ is a proof, so
its type is true, so $\sem f\rho$ takes equal values, as sets, on
$R$-related elements and hence on the members of a class.
\emph{Elimination of equality.} If $e:a=_Ax$ in the context of the rule
then the truth value of $a=_Ax$ is $\top$, that is
$\sem a\rho=\sem x\rho$; so the check succeeds and
$C[a,\refl]$ and $C[x,e]$ have the same code, proofs being all $\star$.
\emph{Proof irrelevance, $\mathsf{propext}$, $\mathsf{em}$.} Every proof is
$\star$; two equivalent propositions have the same truth value, which
is equality of codes; and $\El(\code\Prop)$ has two elements.
\emph{Cumulativity.} $U_i\subseteq U_{i+1}$, truth values are in $U_0$,
and $\code{\U_i}\in U_{i+1}$ because $\lambda_i\cdot\omega<\lambda_{i+1}$.
\end{proof}

\begin{theorem}\label{thm:main}
Assume Conjecture~\ref{conj:rank}. Then the evaluation relative to $\A$ is
a model of $\CIC_\omega+\EM$ in which $\bot$ is empty. So if $\ZF$ proves
Conjecture~\ref{conj:rank} then $\ZF\vdash\Con(\CIC_\omega+\EM)$.
\end{theorem}
\begin{proof}
Lemmas~\ref{lem:closedex} and \ref{lem:sound}; $\sem\bot\rho$ is the
truth value of $\|\mathbf 0\|$, which is $\bot$, and $\El(\bot)=\emptyset$.
\end{proof}

Lemma~\ref{lem:sound} is written at the level of its cases and has not
been checked against every rule of Section~\ref{sec:tt}. Three points
deserve a second look when it is: that the sorting and the proof marks
which the evaluation reads are determined by the derivation and are
stable under conversion; that the accessible part computed in the
decoding agrees with the inductive reading of $\Acc$ that the
introduction rule needs; and the substitution lemma under sanitized
families.

\section{The rank lemma}\label{sec:ranklemma}

Everything rests on Conjecture~\ref{conj:rank}. We say why we believe
it, what a proof has to contain, and what is known to obstruct the
obvious proof.

\paragraph{The first-order case.}
For a term all of whose variables are free and of small type the
conjecture is the rank lemma of \cite[\S3]{paper2} in type-theoretic
dress. A type former adds a bounded amount to the ranks of its
components. $\ind^\N$ with a type-valued motive iterates the bound of its
step, and the supremum over $n<\omega$ of the iterates is an ordinal. A
$W$-recursion over a tree with small branching iterates the bound of its
step along the height of the tree, which is below the Hartogs number of
its branching types. An $\Acc$-recursion needs an idea: the accessible
part of a relation on a large carrier has unbounded height, but the
value at $x$ depends on the recursive function only at the predecessors
at which the step applies it, which are the values of finitely many
subterms in environments varying in bound variables of small type; the
tree of inspected predecessors has small branching, so its height is
below a Hartogs number, as for well-founded recursion in \cite{paper2}.
None of these bounds mentions the cut.

\paragraph{Higher types.}
A term may bind a variable of a type that is not small, such as
$G:\Type_N\to\Type_N$, and apply a defined functional to the identity,
whose content is the cut. So the bound for a term of function type must
be a function of the bounds of its arguments, hereditarily: a
\emph{profile}, a majorant in Howard's sense \cite{howard73} with
ordinals for numbers. The statement to prove is a fundamental lemma:
the value of a term has the profile computed from the profiles of its
variables by recursion on the term, and the profiles so computed do not
depend on the cut. Since terms are now typed, and every value lives in
the decoding of a code, profiles can be assigned by recursion on the
code, which is well founded; no universal domain of profiles is needed.
Three things are known to be delicate.
\begin{itemize}
\item \emph{Bounds are monotone but not continuous}, because the rank of
  $\Sigma n{:}\N.\,B(n)$ is a genuine supremum. So a recursion cannot be
  bounded by a least fixed point: the step
  $\mathrm{rec}\mapsto(\Sigma n{:}\N.\,\mathrm{rec}\,(p\,x\,n))\to\Prop$ has the
  bound $o\mapsto o+k$, whose only fixed point is ``no bound'', while the
  truth is a countable ordinal. The inspection argument is essential.
\item \emph{The inspected predecessors are not simply those at which the
  recursive function is applied.} Let $K$ be a parameter of bounded
  content. A step of the form
  $K(\lambda X{:}\Type_N.\,\mathrm{rec}\,(h\,X))$ evaluates the recursive
  function at $h(X)$ for every $X$ in the universe.
  What is small is the set of predecessors whose values can reach the
  \emph{rank} of the result, the value of $K$ being bounded whatever it
  is given. So the profile of the recursive function must depend on the
  actual predecessor and not only on a bound for it.
\item \emph{Large elimination makes the shape of a value depend on
  data.} For $B:\N\to\Type_{N+1}$ with $B\,0=\Type_N$ and
  $B\,(n{+}1)=B\,n\to B\,n$, an element of $B\,n$ is a functional of order
  $n$ over the universe. Profiles indexed by codes handle this, but they
  must be joined fibrewise and never across the branches of an
  eliminator whose branches have different levels
  (Example~\ref{ex:hetero}).
\end{itemize}

\paragraph{Why the cut should not matter.}
The cut enters the evaluation in four ways. It is the range of a bound
variable of universe type; but a type former whose domain is not in
$U_N$ has a code of rank at least the cut, which the conjecture does not
constrain, while one whose domain is in $U_N$ binds a variable whose
values have content bounded by the rank of the domain; and an
abstraction over a larger type is only ever applied to arguments that
have themselves been bounded. It is the threshold at which families are
sanitized; but a fibre replaced by $\bot$ only lowers the rank. It
decides when an application is junk, because the argument fell outside
the universe; but junk is $\emptyset$. And it changes the truth of
propositions, through the oracle and through the extents; but the
conjecture quantifies over all oracles and extents, and a proposition
can influence a type only through the relation of a quotient, the
domain of an $\Acc$-recursion, the check in $\ind^=$, and the decoding of
a proof type, each of which selects among alternatives that are bounded
anyway.

\paragraph{Attacks that failed.}
In the extensional setting of \cite[\S8]{paper2}, where the conjecture
was stated for untyped terms and arbitrary truth assignments and is
therefore stronger, the following attempts to produce types just below
the cut from parameters of bounded content failed: ill-typed terms
binding a universe; parameters of content $0$ over a large carrier
without codes, such as a well-order of length above the cut on
$\Type_N\to\Prop$ with a countable predecessor function, driving an
$\Acc$-recursion of a Hartogs-raising step (the inspected tree is
countable, and after an overflow the junk value restarts the iteration
from the bottom); trees of large height and content $0$ with large
branching; classes of a quotient of the universe containing both a
small code and a code near the cut; and selection of a code by a
proposition, which would need a description operator. The present
formulation has not yet been attacked on its own terms, and the oracle
is a new degree of freedom for an adversary: an oracle may contain a
cofinal family, and the conjecture claims that a term with bounded
parameters cannot get at it.

\section{Examples}\label{sec:examples}

\begin{example}[a climbing function]\label{ex:climb}
Suppose $\lambda_0$ has cofinality $\omega$; fix a cofinal sequence
$(\beta_n)$, let $\mathrm{next}(\beta)$ be the least $\beta_n>\beta$, and
let $F(X)$ be a code of rank $\mathrm{next}(\rk X)$. Every pointwise
condition on $F$, such as a Howard majorant mapping ordinals below
$\lambda_0$ to ordinals below $\lambda_0$, accepts $F$; yet
$\Sigma n{:}\N.\,F^n(\N):\Type_0$ denotes a code of rank $\lambda_0$. The
extent of a universe below an inaccessible is singular, so some small
index type always admits a cofinal family, and no condition on values
excludes it. In the model $F$ is excluded because it is not generated:
if $F=\sem s^\X\rho$ with $\rho$ of content below $\gamma<\lambda_0$ then
the term $\ind^\N\,\N\,(\lambda\_.\,s)\,x$, with $x{:}\N$ free, takes the
values $F^n(\N)$, which lie in $U_0$ and are cofinal in $\lambda_0$,
contradicting closedness.
\end{example}

\begin{example}[trees and Hartogs numbers]\label{ex:set}
Let $\mathsf{set}:=Wx{:}\Type_0.\,x$ and define $T:\mathsf{set}\to\Type_0$ by
$W$-recursion, $T(\mathsf{sup}\,X\,d):=(\Sigma x{:}X.\,T(d\,x))\to\Prop$. Then
$\rk T(w)$ is at least the height of $w$. A tree in the decoding under
the oracle $V_\vartheta$ over a label of rank below $\gamma$ can have any
height below a Hartogs number of $V_{\gamma\cdot\omega}$, and $T(w)$ with
$w$ free of content below $\gamma$ is then a type of that rank; so an
$N$-closed ordinal is closed under $\beta\mapsto\hartogs(V_{\beta+1})$, and
in $\ZFC$ it is a strong limit cardinal. In the model the subtree
functions are generated, and the heights of the trees of the model are
bounded below $\lambda_0$ for that reason.
\end{example}

\begin{example}[heterogeneous fibres]\label{ex:hetero}
Let $B'(0):=\Type_5$ and $B'(1):=\Type_0$. A family
$f:\N\to\Pi b{:}\mathbf 2.\,B'(b)$ can have ranks at $1$ cofinal in
$\lambda_0$ while being dominated, in any order that joins the fibres of
a dependent function, by a constant family that is large at $0$. Content
is computed fibre by fibre for this reason, and so must profiles be.
\end{example}

\begin{example}[why closedness is stated in the structure cut at $\lambda$ itself]\label{ex:sat}
A tempting alternative defines $\lambda_N$ as the least ordinal closed
under functions $\nu_t(\alpha)$ computed in the structure cut at
$\alpha<\lambda$, and compares the model with those structures
afterwards. It fails. In a structure cut at $\alpha$ the recursive
function of a recursion, as a bound variable of the step, ranges over
functions into the cut universe, so the computed bound of any
rank-raising recursion saturates near $\alpha$, every $\nu_t(\alpha)$ is
about $\alpha$, and the least closed ordinal is just the least ordinal
$\eta$ closed under $\alpha\mapsto\hartogs(V_{\alpha+1})$. But for
$X:\Type_N$ and a well-founded linear order $o=(R,\dots)$ on $X$ define
$T_o:X\to\Type_N$ by $\Acc$-recursion,
$T_o(x):=(\Sigma y{:}X.\,\Sigma h{:}R\,y\,x.\,T_o(y))\to\Prop$, and put
$t(X):=\Sigma o.\,\Sigma x{:}X.\,T_o(x)$. Then $\rk t(X)\ge\hartogs(\El(X))$,
the decoding of $T_o(x)$ is a cumulative hierarchy of height the order
type of $x$, the ranks of $t^n(\N)$ are cofinal in $\eta$, and
$\Sigma n{:}\N.\,t^n(\N):\Type_N$ has rank $\eta$. Under
Definition~\ref{def:closedord} the values $t^n(\N)$ are honest in the
structure cut at $\lambda$ as long as they lie below $\lambda$, so
$\lambda_N>\eta$; and no term can diagonalize against $\lambda_N$,
because the sequence in Lemma~\ref{lem:closedex} uses the bounds of all
terms at once.
\end{example}

\begin{example}[why the eliminator of equality checks the equation]\label{ex:confusion}
Closedness quantifies over oracles and environments that are not those
of the model, and at such an environment a proposition may be true for
no good reason. Let $B$ be a parameter with $B(f)=\code{\U_{N+1}}$ and
$B(g)=\code{\U_N}$ for two elements $f,g$, in a context with
$e:B\,f=B\,g$ and $Y:B\,f$. If $\ind^=$ evaluated to the identity
whenever a proof of the equation is supplied, the transport of $Y$ along
$e$ would be a term of type $\Type_N$ whose value is an arbitrary code of
$U_{N+1}$, of rank unbounded in the higher extents, and no ordinal would
be closed. Here the truth value of $B\,f=B\,g$ is computed, it is
$\bot$, and the context has no environment at all; and in general the
check makes the value junk unless the two sides are equal. In the
extensional theory the same confusion arises from reflection with no
eliminator to guard, which is one reason this paper is about the
intensional theory. Sanitized families answer the remaining case, a
well-typed family that takes a value outside its universe at an
argument outside the model.
\end{example}

\section{Status}\label{sec:status}

\emph{Proved, given the definitions:} the existence of closed ordinals
from the rank lemma (Lemma~\ref{lem:closedex}); the closure properties of
the generated elements (Lemma~\ref{lem:gen}); type formation
(Lemma~\ref{lem:tyform}). \emph{Proved by cases, to be checked rule by
rule:} soundness (Lemma~\ref{lem:sound}). Known loose ends there: that
the class of a quotient of a large type is generated, which
Lemma~\ref{lem:gen}(4) obtains from the relation as a parameter of
content $0$ and which needs the relation's curried function to be
admitted by an oracle without disturbing the decoding of the carrier;
and the three points listed after Theorem~\ref{thm:main}.
\emph{Open:} Conjecture~\ref{conj:rank}, on which the theorem depends.
\emph{Not attempted:} general inductive families; the refinement to
finitely many universes, $\ZFh_{n+1}\vdash\Con(\CIC_{n+2}+\EM)$, which
would close the hierarchy of \cite{paper2}.

This paper supersedes \cite[\S8]{paper2}. That section went through
three designs: Howard-style majorants with closure functions computed in
cut structures and a retraction between structures, which fails by
Example~\ref{ex:sat}; a partial equivalence relation over full function
spaces for the extensional theory, with typed closedness under arbitrary
truth assignments, which is unsatisfiable by the extensional form of
Example~\ref{ex:confusion}; and the same with closedness as a gap
condition on untyped terms, for which the rank lemma is plausible but
resists proof because untyped terms include junk self-applications that
are typable under reflection. The present design keeps from them the
codes, the notion of content, the gap form of closedness, and the
examples, and replaces truth assignments and the partial equivalence
relation by the oracle.

{\small
\bibliographystyle{plain}
\bibliography{references}}

\end{document}
